\section{SCA算法收敛性分析}

本节基于第3章提出的逐次凸逼近（SCA）轨迹优化算法，通过数值实验验证其收敛特性的三个核心方面：单阶段内的迭代收敛行为、跨阶段的目标函数演化趋势，以及CRB与通信速率在收敛过程中的耦合关系。实验设置权衡系数$\eta=0.7$，每阶段航点数$N_m=20$，SCA最大迭代次数$I_{\max}=50$。实验产出的两条UAV轨迹共支撑了3个优化阶段的SCA迭代，累计66轮迭代。

\subsection{全局收敛趋势}

图\ref{fig:s1_global}以连续累积迭代次数为横轴，展示了目标函数$J$、CRB以及平均通信速率$\bar{R}_{\text{com}}$跨三个阶段的总体演化趋势。表\ref{tab:sca_stages}汇总了各阶段首末迭代的关键指标数值。

\begin{table}[htbp]
\centering
\caption{SCA各阶段收敛数据汇总}
\label{tab:sca_stages}
\begin{tabular}{ccccc}
\toprule
阶段 & 迭代次数 & 目标函数$J$（初$\rightarrow$终） & CRB（初$\rightarrow$终） & 速率$\bar{R}_{\text{com}}$（初$\rightarrow$终） \\
\midrule
Stage 1 & 17 & $12228\rightarrow1101$ & $1.79\times10^{4}\rightarrow2.03\times10^{3}$ & $9.58\rightarrow10.62$ Mbps \\
Stage 2 & 16 & $137\rightarrow-188$ & $6.00\times10^{2}\rightarrow1.23\times10^{2}$ & $9.45\rightarrow9.15$ Mbps \\
Stage 3 & 33 & $-223\rightarrow-241$ & $3.38\times10^{1}\rightarrow4.06$ & $8.23\rightarrow8.14$ Mbps \\
\bottomrule
\end{tabular}
\end{table}

从图\ref{fig:s1_global}和表\ref{tab:sca_stages}中可以得出以下核心观察：

\textbf{（1）CRB的阶梯式数量级下降。}
CRB在三个阶段间呈现显著的阶梯式下降特征——从Stage~1初始的$1.79\times10^{4}$降至Stage~3最终的$4.06$，累计降幅达\textbf{三个数量级以上}（约$4.4\times10^{3}$倍）。每一阶段的CRB下降对应着不同的物理机制：Stage~1的CRB骤降源于UAV从粗扫描的4个等高度悬停点过渡至第一个三维优化轨迹，高度自由度的释放使FIM条件数发生质变；Stage~2的CRB改善得益于目标估计更新后UAV轨迹更逼近目标真实位置；Stage~3的CRB精调则是在目标估计已趋于收敛（位置误差降至约10~m量级）的条件下，对悬停点空间布局的进一步微调。

\textbf{（2）目标函数的符号翻转。}
目标函数$J=\eta\cdot\text{CRB}-(1-\eta)\bar{R}_{\text{com}}/\kappa$在Stage~1内保持正值（$12228\rightarrow1101$），表明CRB项主导优化方向；在Stage~2初期穿越零线（$137\rightarrow-188$），标志着CRB已降至与归一化速率项相比拟的量级，通信性能开始在目标函数中体现权重；在Stage~3中$J$始终为负并趋于平稳（$-223\rightarrow-241$，变化幅度仅$8.1\%$），说明系统已进入局部最优邻域。目标函数符号的翻转并非算法异常，而是多目标优化中CRB从主导项退化为均衡项的必然结果，反映了$\eta=0.7$设置下通信与感知之间的自然均衡点。

\textbf{（3）通信速率的有限牺牲。}
三阶段间平均通信速率从$9.58$~Mbps缓降至$8.14$~Mbps，累计降幅约$15\%$。值得注意的是速率的下降并非单调发生在每一阶段内部——Stage~1的SCA迭代中速率反而从$9.58$提升至$10.62$~Mbps（增幅$10.9\%$），这是因为初始直线轨迹向用户-目标连线中点方向优化时，UAV同时受益于向用户方向的偏移和向目标方向的偏移。Stage~2和Stage~3的轻微速率下降（分别降$3.2\%$和$1.1\%$）则是UAV为进一步压低CRB而向目标方向微调轨迹所付出的边际代价。

\subsection{单阶段内迭代收敛特性}

图\ref{fig:s1_perstage}以分阶段网格图的形式展示了每阶段内部目标函数$J$、CRB及通信速率随SCA迭代次数的演化曲线。该图揭示了SCA迭代的三个关键行为特征：

\textbf{（1）前5轮快速收敛。}
在三个阶段的SCA迭代中，目标函数和CRB均在前5轮迭代内完成绝大部分改善。以Stage~1为例：目标函数在第1至第5轮迭代中从$12228$骤降至约$1585$（降幅$87\%$），此后12轮迭代仅进一步降至$1101$。这一特征与SCA的理论性质一致——凸近似在初始点附近与原问题的梯度最为吻合，每轮迭代沿近似梯度方向取得最大下降；随着迭代点逼近局部最优，凸近似与原问题之间的偏差逐渐暴露，收敛速率从线性退化为亚线性。

\textbf{（2）线搜索对振荡的抑制。}
从图\ref{fig:s1_perstage}中可观察到，尽管SCA每轮求解的是原问题的凸近似，但目标函数$J$、CRB及速率在所有三个阶段中均呈现\textbf{严格的单调变化}——$J$单调下降，CRB单调下降（Stage~1速率除外），未出现任何振荡或发散现象。这验证了第3章设计的线搜索策略（候选步长集$\{1.0,0.8,\dots,0.05\}$）的有效性：当凸近似给出的搜索方向在原始非凸目标函数上不能保证下降时，线搜索通过缩减步长寻找可行的下降点，避免了SCA常见的“近似-原始不一致”导致的振荡问题。

\textbf{（3）Stage~3收敛减缓的机理。}
Stage~3耗费了33轮迭代才收敛，显著多于Stage~1（17轮）和Stage~2（16轮）。这一现象的深层原因有三：（a）Stage~3初始CRB已低至$33.79$，进一步改善的相对空间缩小，目标函数梯度在局部最优附近趋于平坦；（b）此时剩余能量预算紧张（约$13$~kJ），UAV轨迹在更紧的能量约束下进行微调，可行域缩小限制了每轮迭代的改善幅度；（c）Stage~3的目标函数绝对值较小（$-223$至$-241$），在收敛容差$\varepsilon_{\text{obj}}=10^{-3}$下，相对变化需小于$10^{-5}$量级才能触发终止条件，因而需要更多迭代来跨越这一平坦区域。尽管如此，33轮迭代后算法仍自然收敛（未触及$I_{\max}=50$上限），表明所设收敛容差和最大迭代次数是合理的。

\subsection{CRB与速率的耦合-解耦分析}

图\ref{fig:s1_global}和\ref{fig:s1_perstage}共同揭示了CRB与通信速率之间有趣的“先耦合、后解耦”关系：

在Stage~1的优化初期，UAV轨迹从远离目标和用户的初始位置出发，CRB的剧烈下降与速率的同步提升呈现\textbf{正耦合}——两者同时改善。这是因为初始轨迹的任何优化方向（无论是偏向用户还是偏向目标）相对于初始的直线轨迹均是改进，尚未触及通信与感知之间的本质矛盾。

随着优化深入（Stage~2及Stage~3），CRB与速率的关系转变为\textbf{负耦合}（权衡关系）——进一步降低CRB需要以牺牲速率为代价。这一转变标志着UAV轨迹已进入用户与目标之间的“博弈区”，任何空间位移对一方的增益必然伴随着对另一方的损失，是ISAC系统中通信-感知内在权衡在几何层面的直接体现。

\subsection{收敛性与算法效能小结}

综合以上分析，SCA轨迹优化算法在数值实验中表现出良好的收敛特性：目标函数单调下降且无振荡，三个阶段的CRB累计改善超过三个数量级，单阶段求解时间约$4.8$~s（MOSEK求解器），66轮总SCA迭代在约$14.4$~s内完成。这一计算效率完全满足离线轨迹规划的实时性要求。值得注意的是，最终CRB（$4.06$）意味着目标定位的理论精度下界约为$\sqrt{4.06}\approx2.0$~m，为后续定位算法（第4章）提供了高质量的观测几何基础。
